\documentclass[aspectratio=169]{beamer}

\usetheme{Madrid}
\usecolortheme{default}
\setbeamertemplate{navigation symbols}{}

\usepackage[T1]{fontenc}
\usepackage[utf8]{inputenc}
\usepackage{lmodern}
\usepackage{booktabs}
\usepackage{array}

\definecolor{brand}{RGB}{15,118,110}
\definecolor{branddark}{RGB}{11,94,88}
\definecolor{accent}{RGB}{3,105,161}
\definecolor{textgray}{RGB}{70,84,97}

\setbeamercolor{structure}{fg=brand}
\setbeamercolor{title}{fg=branddark}
\setbeamercolor{frametitle}{fg=branddark}
\setbeamercolor{normal text}{fg=black,bg=white}
\setbeamercolor{block title}{fg=white,bg=brand}
\setbeamercolor{block body}{fg=black,bg=brand!6}
\setbeamercolor{alerted text}{fg=accent}

\title{Web Streaming App}
\subtitle{Janus VideoRoom + WebRTC Streaming MVP}
\author{Daniela}
\date{\today}

\begin{document}

\begin{frame}
  \titlepage
\end{frame}

\begin{frame}{Presentation Goal}
  \begin{itemize}
    \item Present a working streaming MVP built on \textbf{WebRTC} and \textbf{Janus VideoRoom}
    \item Show that the solution is not only functional, but also \textbf{structured, resilient, and extensible}
    \item Explain the project in a way that is easy to follow for both technical review and project fit discussion
  \end{itemize}
\end{frame}

\begin{frame}{What The Application Does}
  \begin{block}{Core Use Case}
    A user can start the app either as a \textbf{Publisher} to stream camera and microphone, or as a \textbf{Viewer} to watch the live stream.
  \end{block}

  \vspace{0.4em}
  \begin{itemize}
    \item Browser client connects directly to Janus over \textbf{raw WebSocket signaling}
    \item Publisher creates or reuses a room and pushes live media
    \item Viewer automatically detects the active publisher feed and subscribes to it
    \item The UI exposes status, logs, friendly errors, and recovery feedback
  \end{itemize}
\end{frame}

\begin{frame}{Why This Project Is Valuable}
  \begin{itemize}
    \item Demonstrates \textbf{real-time communication} with low-latency media delivery
    \item Uses a lightweight architecture: no heavy frontend framework, no Janus helper library
    \item Shows understanding of both \textbf{signaling flow} and \textbf{media negotiation}
    \item Includes practical engineering concerns:
    \begin{itemize}
      \item session lifecycle
      \item room management
      \item reconnection handling
      \item user experience during failures
    \end{itemize}
    \item Adds a Java viewer skeleton, which shows the design can evolve beyond a browser-only MVP
  \end{itemize}
\end{frame}

\begin{frame}{High-Level Architecture}
  \begin{center}
    \renewcommand{\arraystretch}{1.4}
    \begin{tabular}{>{\bfseries}m{0.24\linewidth} m{0.68\linewidth}}
      Browser UI & Inputs, buttons, local/remote video elements, status badge, log panel \\
      \addlinespace
      app.js & Janus signaling, WebRTC peer connection, publisher/viewer flows, recovery logic \\
      \addlinespace
      Janus Server & Session management, VideoRoom plugin, signaling transport, media routing \\
      \addlinespace
      Java Client & Subscriber skeleton for external integrations or non-browser consumers \\
    \end{tabular}
  \end{center}

  \vspace{0.6em}
  \begin{block}{Design Choice}
    The browser talks to Janus using the \textbf{janus-protocol} directly, which keeps the implementation transparent and easier to reason about during debugging.
  \end{block}
\end{frame}

\begin{frame}{Publisher Flow}
  \begin{enumerate}
    \item Open WebSocket connection to Janus
    \item Create Janus session
    \item Attach to \texttt{janus.plugin.videoroom}
    \item Check whether the room exists; create it if needed
    \item Ensure there is no active publisher already streaming
    \item Join the room as publisher
    \item Capture camera and microphone with \texttt{getUserMedia}
    \item Create WebRTC offer and send \texttt{configure} request with JSEP
    \item Receive SDP answer from Janus and start publishing
  \end{enumerate}
\end{frame}

\begin{frame}{Viewer Flow}
  \begin{enumerate}
    \item Open WebSocket connection, create session, and attach to VideoRoom
    \item Ensure the room exists
    \item Poll room participants to detect a live publisher
    \item Join as subscriber using the publisher \texttt{feed} and \texttt{streams} mapping
    \item Receive JSEP offer from Janus
    \item Create local SDP answer and send \texttt{start}
    \item Attach the remote stream to the viewer video element
    \item Monitor publisher changes and rejoin automatically if the feed disappears or changes
  \end{enumerate}
\end{frame}

\begin{frame}{What Is Strong In The Implementation}
  \begin{itemize}
    \item \textbf{Clear state model}: idle, connecting, connected, publishing, viewer waiting, reconnecting, error
    \item \textbf{Friendly error mapping}: device permissions, invalid room, missing feed, server availability
    \item \textbf{Automatic recovery}: detects disconnects, resets media state, and guides the user to restart
    \item \textbf{Viewer resiliency}: auto-join and monitor loops help keep playback aligned with the live publisher
    \item \textbf{Good MVP boundaries}: small codebase, readable flow, focused responsibilities
  \end{itemize}
\end{frame}

\begin{frame}{Java Client Extension}
  \begin{block}{Why It Matters}
    The project is not limited to a browser demo. It already includes a Java subscriber skeleton, which is important for backend tooling, monitoring, kiosk playback, or future service integrations.
  \end{block}

  \vspace{0.4em}
  \begin{itemize}
    \item Mirrors the Janus lifecycle: create session, attach plugin, keepalive, subscribe, start, trickle ICE
    \item Uses a dedicated WebRTC engine abstraction
    \item Waits for Janus asynchronous offer events and responds with a local SDP answer
    \item Shows that the signaling model is reusable across platforms
  \end{itemize}
\end{frame}

\begin{frame}{Key Technical Decisions}
  \begin{itemize}
    \item \textbf{Raw Janus WebSocket API instead of janus.js}
    \begin{itemize}
      \item better protocol visibility
      \item easier to explain exact signaling steps
      \item less hidden behavior
    \end{itemize}
    \item \textbf{Single active publisher per room}
    \begin{itemize}
      \item simplifies the MVP
      \item avoids ambiguous viewer selection
    \end{itemize}
    \item \textbf{Auto room creation}
    \begin{itemize}
      \item lowers setup friction in development
      \item makes the demo smoother
    \end{itemize}
  \end{itemize}
\end{frame}

\begin{frame}{Current Limitations}
  \begin{itemize}
    \item TURN/STUN configuration is intentionally minimal and suited mainly for local or controlled environments
    \item Security and authentication are not the focus of this MVP
    \item The viewer currently selects the first live publisher, which is enough for a single-publisher room but not for multi-stream discovery
    \item The frontend is plain JavaScript, which is great for clarity, but would need stronger modularization as the project grows
  \end{itemize}
\end{frame}

\begin{frame}{How I Would Evolve It}
  \begin{itemize}
    \item Add authentication and room-level access control
    \item Externalize configuration for Janus endpoints, ICE servers, and room policies
    \item Improve observability with structured logs and metrics
    \item Support multiple publishers and explicit viewer selection
    \item Add automated tests around signaling state transitions and error scenarios
    \item Package the Java client as a reusable integration module
  \end{itemize}
\end{frame}

\begin{frame}{Why This Fits A Real Project}
  \begin{itemize}
    \item The MVP already covers the hard part: \textbf{real-time media + signaling orchestration}
    \item The code shows practical thinking, not only a demo mindset
    \item It is understandable, extendable, and close to production concerns
    \item It creates a solid base for features such as remote support, live monitoring, streaming dashboards, or communication platforms
  \end{itemize}
\end{frame}

\begin{frame}{Closing}
  \begin{block}{Takeaway}
    This project demonstrates the ability to build a clean real-time streaming flow end to end: UI, signaling, WebRTC negotiation, recovery behavior, and extension into Java.
  \end{block}

  \vspace{0.5em}
  \begin{center}
    \Large Questions
  \end{center}
\end{frame}

\end{document}
